\documentclass[12pt,a4paper]{report}

\usepackage[T1]{fontenc}
\usepackage[utf8]{inputenc}
\usepackage[french]{babel}
\usepackage[a4paper,margin=2.5cm]{geometry}
\usepackage{setspace}
\usepackage{hyperref}
\usepackage{enumitem}
\usepackage{graphicx}
\usepackage{float}

\onehalfspacing
\setlist[itemize]{itemsep=4pt,topsep=4pt}
\hypersetup{colorlinks=true,linkcolor=black,urlcolor=blue,citecolor=black}

\begin{document}

% -------------------------
% Page de garde
% -------------------------
\begin{titlepage}
  \centering
  {\large Republique Tunisienne\par}
  {\large Ministere de l'Enseignement Superieur et de la Recherche Scientifique\par}
  \vspace{0.8cm}
  {\Large Nom de l'etablissement\par}
  \vspace{1.2cm}

  {\Large \textbf{Rapport de Projet / PFE}\par}
  \vspace{0.5cm}
  {\large Specialite : \ldots\par}
  \vspace{1.2cm}

  {\Huge \textbf{Deploiement CI/CD d'une application cloud}\par}
  \vspace{0.3cm}
  {\Large Approche PRISM et fiabilisation du pipeline\par}
  \vspace{1.5cm}

  \begin{flushleft}
    \textbf{Realise par :} \ldots\\
    \textbf{Encadre par :} \ldots\\
    \textbf{Classe :} \ldots\\
    \textbf{Annee universitaire :} 2025--2026
  \end{flushleft}

  \vfill
  {\large Date : \ldots\par}
\end{titlepage}

% -------------------------
% Dedicace et remerciements
% -------------------------
\chapter*{Dedicace}
\addcontentsline{toc}{chapter}{Dedicace}
Je dedie ce travail a ma famille, a mes enseignants et a toute personne ayant contribue a mon apprentissage.

\chapter*{Remerciements}
\addcontentsline{toc}{chapter}{Remerciements}
Je remercie mon encadrant, l'equipe pedagogique et mes collegues pour leur accompagnement,
leurs conseils et leur soutien tout au long de ce projet.

% -------------------------
% Front matter
% -------------------------
\chapter*{Resume}
\addcontentsline{toc}{chapter}{Resume}
Ce rapport presente une demarche de deploiement continu sur infrastructure cloud,
avec une approche PRISM orientee robustesse, observabilite et automatisation.
Le travail met l'accent sur la gestion des incidents et la validation de la mise en production.

\chapter*{Abstract}
\addcontentsline{toc}{chapter}{Abstract}
This report presents a CI/CD deployment approach on cloud infrastructure,
using a PRISM framework focused on reliability, observability, and automation.
The work highlights incident handling and production-readiness validation.

\tableofcontents
\listoffigures
\listoftables
\clearpage

\chapter*{Liste des acronymes}
\addcontentsline{toc}{chapter}{Liste des acronymes}
\begin{itemize}
  \item CI/CD : Continuous Integration / Continuous Deployment
  \item EC2 : Elastic Compute Cloud
  \item GHCR : GitHub Container Registry
  \item IaC : Infrastructure as Code
  \item SSH : Secure Shell
  \item SLA : Service Level Agreement
\end{itemize}

\chapter*{Presentation}
\addcontentsline{toc}{chapter}{Presentation}
Ce document presente une demarche de deploiement continue d'une application web sur infrastructure cloud.
L'objectif est pedagogique : expliquer une approche CI/CD fiable, reproductible et maintenable,
avec un accent sur la gestion des incidents et la validation finale du service.

Ce rapport est concu pour rester dans une plage de 35 a 40 pages en version finale,
avec un niveau de detail oriente deploiement et non fonctionnel metier.

\chapter{Contexte du Projet}
\section{Objectif general}
Le projet vise a automatiser la chaine complete de livraison :
\begin{itemize}
  \item construction des images,
  \item publication dans un registre,
  \item provisionning infrastructure,
  \item deploiement automatise,
  \item verification de sante post-deploiement.
\end{itemize}

\section{Cadre PRISM}
L'approche PRISM structure la methode de travail :

\subsection{P - Preparation}
\begin{itemize}
  \item preparation de l'environnement (cloud, secrets, depots, acces SSH),
  \item preparation des fichiers d'orchestration et de configuration.
\end{itemize}

\subsection{R - Realisation}
\begin{itemize}
  \item construction automatique des images,
  \item publication vers un registre de conteneurs.
\end{itemize}

\subsection{I - Integration}
\begin{itemize}
  \item provisionning de l'infrastructure,
  \item connexion du pipeline avec la machine cible.
\end{itemize}

\subsection{S - Stabilisation}
\begin{itemize}
  \item verification de disponibilite SSH,
  \item relances controlees, timeouts et fallback,
  \item verification interne de sante des services.
\end{itemize}

\subsection{M - Monitoring / Maintenance}
\begin{itemize}
  \item diagnostics automatiques en cas d'echec,
  \item ameliorations incrementales du workflow.
\end{itemize}

\chapter{Architecture}
\section{Vue generale de deploiement}
Elements principaux :
\begin{itemize}
  \item pipeline GitHub Actions (jobs build-images, infra, deploy),
  \item instance cloud (EC2),
  \item orchestration Docker Compose,
  \item registre d'images GHCR,
  \item verification de sante interne et externe.
\end{itemize}

Resultat attendu : une nouvelle version est construite, publiee, deployee et verifiee automatiquement.

\begin{figure}[H]
  \centering
  \fbox{\parbox[c][5cm][c]{0.85\textwidth}{\centering
  Capture architecture globale (a remplacer)}}
  \caption{Architecture generale de deploiement}
\end{figure}

\section{Flux d'execution}
\begin{enumerate}
  \item Build des images et publication.
  \item Application de l'infrastructure.
  \item Verification SSH et readiness de la cible.
  \item Deploiement compose et controles de sante.
  \item Publication URL finale et diagnostics en cas d'incident.
\end{enumerate}

\chapter{Partie Conception}
\section{Conception du pipeline CI/CD}
Le workflow est concu avec trois blocs logiques :
\begin{itemize}
  \item \textbf{Build} : generation et push des images backend/frontend.
  \item \textbf{Infra} : Terraform init/apply et extraction des sorties utiles.
  \item \textbf{Deploy} : execution distante securisee, orchestration et verifications.
\end{itemize}

\section{Conception de la robustesse}
Les mecanismes de fiabilite integres sont :
\begin{itemize}
  \item retries sur les etapes sensibles,
  \item timeouts sur les commandes potentiellement bloquantes,
  \item fallback d'infrastructure en cas d'instance non exploitable,
  \item fallback de build en cas d'echec Buildx,
  \item collecte de diagnostics en fin d'echec.
\end{itemize}

\section{Conception des criteres de validation}
Le deploiement est valide si :
\begin{itemize}
  \item images construites et publiees,
  \item cible joignable en SSH,
  \item services en etat Up/Healthy,
  \item endpoint de health en HTTP 200,
  \item URL finale accessible.
\end{itemize}

\chapter{Partie Realisation}
\section{Scenario 1 : Build et publication des images}
\textbf{Objectif :} produire et publier les images versionnees.

\textbf{Etapes :}
\begin{itemize}
  \item setup du runner,
  \item build backend/frontend,
  \item push GHCR avec tag SHA et latest,
  \item activation du chemin fallback si Buildx indisponible.
\end{itemize}

\textbf{Resultat attendu :} images disponibles dans le registre.

\begin{figure}[H]
  \centering
  \fbox{\parbox[c][5cm][c]{0.85\textwidth}{\centering
  Capture logs build et push (a remplacer)}}
  \caption{Scenario 1 - Build et publication}
\end{figure}

\section{Scenario 2 : Provisionning et acces cible}
\textbf{Objectif :} preparer l'environnement cloud pour le deploiement.

\textbf{Etapes :}
\begin{itemize}
  \item Terraform init/apply,
  \item lecture des outputs (IP publique, instance id),
  \item verification des status checks,
  \item verification de readiness SSH avec retries.
\end{itemize}

\textbf{Resultat attendu :} instance operationnelle et joignable.

\begin{figure}[H]
  \centering
  \fbox{\parbox[c][5cm][c]{0.85\textwidth}{\centering
  Capture Terraform et SSH readiness (a remplacer)}}
  \caption{Scenario 2 - Provisionning et acces}
\end{figure}

\section{Scenario 3 : Deploiement et verification de sante}
\textbf{Objectif :} lancer les services et valider la disponibilite.

\textbf{Etapes :}
\begin{itemize}
  \item pull des images,
  \item lancement Docker Compose,
  \item health checks internes puis externes,
  \item collecte de logs en cas d'echec.
\end{itemize}

\textbf{Resultat attendu :} service accessible avec HTTP 200 sur health.

\begin{figure}[H]
  \centering
  \fbox{\parbox[c][5cm][c]{0.85\textwidth}{\centering
  Capture status conteneurs et health checks (a remplacer)}}
  \caption{Scenario 3 - Deploiement et validation}
\end{figure}

\section{Incidents et corrections appliquees}
\begin{enumerate}
  \item Timeout SSH malgre instance active.
  \item Blocages sur commandes non bornees (status/probes).
  \item Lenteur de readiness backend.
  \item Echec Buildx lie a des timeouts reseau externes.
\end{enumerate}

Pour chaque incident, la correction a suivi la logique : identifier, borner, ajouter fallback, mesurer.

\chapter*{Conclusion Generale}
\addcontentsline{toc}{chapter}{Conclusion Generale}
La demarche PRISM a permis de structurer un deploiement CI/CD plus fiable et plus pedagogique.
La valeur principale est la maitrise d'un processus reproductible, observable et ameliorable,
au-dela de la seule mise en ligne technique.

Perspectives :
\begin{itemize}
  \item ajouter des tests automatiques avant deploiement,
  \item integrer rollback automatique,
  \item renforcer monitoring metriques et alertes,
  \item suivre des indicateurs SLA simples (disponibilite, temps de deploiement).
\end{itemize}

\appendix
\chapter{Annexes}
\section{Captures d'ecran a remplacer}
\begin{itemize}
  \item Capture 1 : vue globale du workflow CI/CD.
  \item Capture 2 : logs build/push GHCR.
  \item Capture 3 : output Terraform + IP cible.
  \item Capture 4 : verification SSH et readiness.
  \item Capture 5 : demarrage Docker Compose.
  \item Capture 6 : health checks internes/externes.
  \item Capture 7 : URL finale accessible.
  \item Capture 8 : logs diagnostics en cas d'echec.
\end{itemize}

\chapter*{Bibliographie}
\addcontentsline{toc}{chapter}{Bibliographie}
\begin{thebibliography}{9}
  \bibitem{gha}
  GitHub, \textit{GitHub Actions Documentation},\
  \url{https://docs.github.com/actions}

  \bibitem{docker}
  Docker, \textit{Docker Documentation},\
  \url{https://docs.docker.com}

  \bibitem{terraform}
  HashiCorp, \textit{Terraform Documentation},\
  \url{https://developer.hashicorp.com/terraform/docs}

  \bibitem{aws}
  Amazon Web Services, \textit{Amazon EC2 Documentation},\
  \url{https://docs.aws.amazon.com/ec2/}
\end{thebibliography}

\end{document}
